
Machine learning can be an attractive tool for either a primary detection capability or supplementary detection heuristics.  Supervised learning models automatically exploit complex relationships between file attributes in training data that are discriminating between malicious and benign samples. 
Furthermore, properly regularized machine learning models generalize to new samples whose features and labels follow a similar distribution to the training data. 

However, malware detection using machine learning has not received nearly the same attention in the open research community as other applications, where rich benchmark datasets exist.   These include handwritten digit classification (e.g., MNIST \cite{MNIST}), image labeling (e.g., CIFAR \cite{CIFAR} or ImageNet \cite{imagenet}), traffic sign detection \cite{trafficsigns}, speech recognition (e.g., TIMIT \cite{TIMIT}), sentiment analysis (e.g., Sentiment140 \cite{sentiment140}), and a host of other datasets suitable for training models to mimic human perception and cognition tasks.  The challenges to releasing a benchmark dataset for malware detection are many, and may include the following.
\begin{itemize}
\item \textit{Legal restrictions.} Malicious binaries are shared generously through sites like VirusShare \cite{virusshare} and VX Heaven \cite{vxheaven}, but benign binaries are often protected by copyright laws that prevent sharing.  Both benign and malicious binaries may be obtained at volume for internal use through for-pay services such as VirusTotal \cite{virustotal}, but subsequent sharing is prohibited.

\item \textit{Labeling challenges.}  Unlike images, text and speech---which may be labeled relatively quickly, and in many cases by a non-expert \cite{buhrmester2011amazon}---determining whether a binary file is malicious or benign can be a time-consuming process for even the well-trained.  The work of labeling may be automated via antimalware scanners that codify much of this human expertise, but the results may be proprietary or otherwise protected.  Aggregating services like VirusTotal specifically restrict the public sharing of vendor antimalware labels \cite{virustotal}.

\item \textit{Security liability and precautions.} There may be risk in promoting a large dataset that includes malicious binaries to a general non-infosec audience not accustomed to taking appropriate precautions such as sandboxed hosting.
\end{itemize}

We address these issues with the release of the Endgame Malware BEnchmark for Research (EMBER) dataset\footnote{Data and code available at \SOURCECODE}, extracted from a large corpus of Windows portable executable (PE) malicious and benign files.  This allows free dissemination of both malicious and benign entities without legal or security concerns.  Samples are released together with the sha256 hash of the original file, and a label to denote whether the file is deemed to be malicious or benign.  The pre-selection of features naturally limits the flexibility of researchers in comparing feature sets.  This is somewhat ameliorated by our release of open source code to compute the PE features for feature comparison studies. The lack of raw binaries also precludes experiments using featureless deep-learning malware detectors (e.g., \cite{raff2017malware}).  However, we hope that by releasing the sha256 hashes, feature extraction source code, as well as a high-performing baseline classifier computed from a subset of the features, this dataset and model codebase will still become a relevant baseline for machine learning malware detection research, and to which featureless deep learning studies may compare.  We demonstrate such a comparison in Section \ref{sec:results}.

We begin in Section \ref{sec:background} with relevant background about the PE file format, as well as a summary of related datasets and approaches for static malware classification.  In Section \ref{sec:method}, we describe the dataset and our methodology for its format.  We demonstrate the efficacy of our baseline model trained on this dataset in Section \ref{sec:results}. Source code and data can be found at \SOURCECODE.
